% =====================================================
% 题型：数列求和与局部等比数列最值填空题 (q_seq_max_ratio.tex)
% 知识板块：数列结构
% 主思维：结构归纳与套用
% 副思维：逆向推导方法、最值模型分析、边界
% 错因标签：区间端点重叠导致的三块比例冲突边界遗漏
% 讲义用途：数列压轴思维
% =====================================================
% 难度：★★★ (难题)
% =====================================================
\newcommand{\qSeqMaxRatioStem}{设实数 \(q\) 满足：存在数列 \(\{a_n\}\)，使得对于任意 \(n \in \mathbb{N}^*\)，均有 \(a_1 + a_2 + \cdots + a_{3n} = n^3 + n\)，且 \(\{a_n\}\) 中有某连续 \(9\) 项 \(a_k, a_{k+1}, \cdots, a_{k+8}\) 是公比为 \(q\) 的等比数列。则 \(q\) 的最大值为 \blank{3cm}。}

\newcommand{\qSeqMaxRatioAnswer}{\(\sqrt[3]{\dfrac{5}{2}}\)}

\newcommand{\qSeqMaxRatioSolution}{%
  \textbf{【解题思路】}\par
  本题为数列与等比数列综合压轴题。第一步，利用前 \(3n\) 项和公式 \(S_{3n} = n^3 + n\)，通过分组作差求出连续 3 项构成的项块和 \(T_n = a_{3n-2} + a_{3n-1} + a_{3n}\)；第二步，分析连续 9 项在等比数列下的几何结构，导出 \(q^3 = \dfrac{T_{j+1}}{T_j}\)；第三步，结合比值单调性与边界冲突排查，锁定公比 \(q\) 的最大值。\par\vspace{0.4em}

  \textbf{【解析计算】}\par
  记前 \(3n\) 项和为 \(S_{3n} = n^3 + n\)，定义项块和 \(T_n = a_{3n-2} + a_{3n-1} + a_{3n}\)。由作差法得：
  \[
    \begin{aligned}
      T_n &= S_{3n} - S_{3n-3} \\[0.3em]
          &= (n^3 + n) - \left[(n-1)^3 + (n-1)\right] \\[0.3em]
          &= 3n^2 - 3n + 2.
    \end{aligned}
  \]
  依次算得前几个项块之和：\(T_1 = 2\)，\(T_2 = 8\)，\(T_3 = 20\)，\(T_4 = 38\)。\par\vspace{0.3em}

  若连续 \(9\) 项构成公比为 \(q\) 的等比数列，则必覆盖至少两个相邻的完整项块 \(T_j\) 与 \(T_{j+1}\)，满足：
  \[
    q^3 = \frac{T_{j+1}}{T_j}.
  \]
  考查比值函数：
  \[
    f(j) = \frac{T_{j+1}}{T_j} = \frac{3j^2 + 3j + 2}{3j^2 - 3j + 2} = 1 + \frac{6j}{3j^2 - 3j + 2}.
  \]
  易知 \(f(j)\) 在 \(j \in \mathbb{N}^*\) 上严格递减，故最大比值应尽可能在最小的 \(j\) 处产生。\par\vspace{0.3em}

  \begin{enumerate}[leftmargin=2em, label=(\arabic*), itemsep=0.3em, topsep=0.3em]
    \item 若 \(j = 1\)：连续 9 项必为 \(a_1, \dots, a_9\)，同时覆盖三个项块 \(T_1, T_2, T_3\)。这要求 \(q^3 = \frac{T_2}{T_1} = 4\) 且 \(q^3 = \frac{T_3}{T_2} = 2.5\)，产生矛盾，舍去；
    \item 若 \(j = 2\)：可取连续 9 项为 \(a_2, \dots, a_{10}\)，仅覆盖 \(T_2, T_3\)，由此只需满足：
      \[
        q^3 = \frac{T_3}{T_2} = \frac{20}{8} = \frac{5}{2} \implies q = \sqrt[3]{\frac{5}{2}}.
      \]
  \end{enumerate}

  故 \(q\) 的最大值为 \(\sqrt[3]{\dfrac{5}{2}}\)。%
}

\newcommand{\qSeqMaxRatioDiagnosis}{%
  这道数列压轴填空题难度极高。许多学生能推导出 \(q^3 = \dfrac{T_{j+1}}{T_j}\) 这一核心关系，并利用单调性发现最大值应该在 \(j=1\) 即 \(\sqrt[3]{4}\) 处产生，但却忽视了“连续 9 项”这一长度限制下，当选取 \(j=1\) 时必须强行涵盖 \(T_3\)，从而引发三块比例冲突的边界问题，从而落入陷阱。%
}

\newcommand{\qSeqMaxRatioTeachingNotes}{%
  本题为数列填空的巅峰之作。在教学讲评中，第一步应指导学生通过“新数列法”（将 3 项合并为一个大项 \(T_n\)）进行作差降维；第二步是重点：不仅要证明最大值在哪，更要通过“存在性验证”（回代具体的项区间如 \(a_2 \sim a_{10}\)）来检验边界条件。这能极大训练学生的逻辑严密性。%
}
